\documentclass{article}
\usepackage[utf8]{inputenc}
\usepackage[T1]{fontenc}
\usepackage{geometry}
\geometry{
    a4paper,
    left=2.5cm,
    right=2.5cm,
    top=2.5cm,
    bottom=2.5cm
}
\usepackage{graphicx}
\usepackage{float}
\usepackage{hyperref}
\usepackage[english]{babel}
\usepackage{amsmath}
\usepackage{amssymb}
\usepackage{array}
\usepackage{booktabs}
\usepackage{caption}
\usepackage{subcaption}
\usepackage{enumitem}
\usepackage{fancyhdr}
\usepackage{titlesec}
\usepackage{url}
\usepackage{textcomp}
\usepackage{xcolor}
\usepackage{listings}
\usepackage{multirow}
\usepackage{longtable}
\usepackage{tabularx}
\usepackage{times}

% Set up page style
\pagestyle{fancy}
\fancyhf{}
\rfoot{\thepage}
\renewcommand{\headrulewidth}{0pt}

% Title formatting
\titleformat{\chapter}[display]{\normalfont\bfseries}{\chaptertitle\thechapter}{0pt}{\Huge}
\titlespacing*{\chapter}{0pt}{50pt}{40pt}

% Document info
\title{Template Document}
\author{Academic Support Agent}
\date{\today}

\begin{document}
\maketitle

\title{\textbf{Title}\\Statement of Research Interests}
\author{Author}
\date{}
\maketitle
This is the place for a summary and an introduction.\\

See Figure \ref{fig:example} for an example Figure with caption.\\

Regarding references: The bibtex file (applications.bib) contains references in an \textit{ad-hoc} format. You can add references and adapt as needed. For example, this is a reference for work in preparation \cite{Doe:inprep}. You can highlight your name with bold font to indicate first author publications \cite{Doe:2022}. This template was used for an application as a postdoctoral researcher in astronomy and astrophysics. You can find a list of journal macros in the preamble. Comment, adapt, and delete as needed.



\section{Research Overview: Add a keyword here}
\label{sec:methods}

\section{Proposed Projects}
\label{sec:proposed_projects}

\subsection{Proposed Project 1}
\label{subsec:proposed_project_1}

\subsubsection{First part of Proposed Project 1}
\subsubsection{Second part of Proposed Project 1}


\subsection{Proposed Project N}
\label{subsec:proposed_project_N}

\begin{figure*}
	% trim={<left> <lower> <right> <upper>}
	\includegraphics[width=\linewidth]{00_Channel_I_horizontal.jpg}
\caption{
Graphical representation of the key stages in the evolution of a massive stellar binary leading to a double neutron star (Figure 1 of Vigna-G\'omez et al. 2020 \cite{2020PASA...37...38V}).
}
\label{fig:example}
\end{figure*}


\begin{multicols}{2}
\setlength{\itemsep}{-10mm}
\end{multicols}
\end{document}